% main.tex — Romanian original
\documentclass[12pt,a4paper]{book}

% Encoding and language
\usepackage[utf8]{inputenc}
\usepackage[T1]{fontenc}
\usepackage[romanian]{babel}

% Maths and layout
\usepackage{amsmath,amssymb,amsthm}
\DeclareMathOperator{\End}{End}
\DeclareMathOperator{\soc}{soc}
\usepackage{geometry}
\geometry{margin=3cm}
\usepackage{graphicx}
\usepackage{microtype}
\usepackage[hidelinks]{hyperref}

% ---------- Theorem environments (adjust later if needed) ----------
% \theoremstyle{plain}
% \newtheorem{theorem}{Teoremă}[chapter]
% \newtheorem{lemma}[theorem]{Lemă}
% \newtheorem{proposition}[theorem]{Propoziţie}
% \newtheorem{corollary}[theorem]{Corolar}

% \theoremstyle{definition}
% \newtheorem{definition}[theorem]{Definiţie}
% \newtheorem{example}[theorem]{Exemplu}

% \theoremstyle{remark}
% \newtheorem{remark}[theorem]{Observaţie}

\newtheorem{theorem}{Teoremă}[section]   % [section] not [chapter]
\newtheorem{lemma}[theorem]{Lemă}
\newtheorem{proposition}[theorem]{Propoziţie}
\newtheorem{corollary}[theorem]{Corolar}
\theoremstyle{definition}
\newtheorem{definition}[theorem]{Definiţie}


% ---------- Title page (fill in the real title later) ----------
\title{DIMENSIUNEA COIREDUCTIBILĂ A INELELOR ŞI MODULELOR}
\author{Coordonator ştiinţific: Prof. Dr. Tiberiu Dumitrescu\\[0.5em]
Student: Sorin Bogde}
\date{Universitatea Bucureşti, Facultatea de Matematică\\[0.5em] 1999}


\begin{document}

\frontmatter
\maketitle
\tableofcontents

% \chapter*{Prefaţă}
% Aici vom introduce eventual un text de prefatare,
% după ce terminăm transcrierea originalului.

\mainmatter

\setcounter{chapter}{-1}
\chapter{Generalităţi}

\section{Module semisimple}

\begin{definition}
Un \(R\)-modul nenul \(S\) se numeşte \emph{simplu} dacă singurele sale submodule sunt
\(0\) şi \(S\).
\end{definition}

\begin{proposition}
Fie \(S\) un \(R\)-modul. Următoarele afirmaţii sunt echivalente:
\begin{enumerate}
  \item \(S\) este modul simplu.
  \item Pentru orice element nenul \(x\in S\) avem \(S = xR\).
  \item \(S \simeq R/I\), unde \(I\) este un ideal drept maximal.
\end{enumerate}
\end{proposition}

\begin{lemma}[Schur]
Fie \(S\) şi \(S'\) două \(R\)-module simple şi \(f\colon S \to S'\) un morfism de \(R\)-module.
Atunci \(f = 0\) sau \(f\) este izomorfism. În particular \(\End_R(S)\) este corp.
\end{lemma}

\begin{definition}
Fie \(M\) un \(R\)-modul şi \((S_i)_{i\in I}\) mulţimea submodulelor simple ale lui \(M\).
Dacă \(M = \sum_{i\in I} S_i\), atunci \(M\) se numeşte \emph{semisimplu}.
\end{definition}

\begin{proposition}
Fie \(M\) un \(R\)-modul semisimplu şi \(N\) un submodul al său.
Atunci există o submulţime \(J \subseteq I\) astfel încât:
\begin{enumerate}
  \item familia \((S_j)_{j\in J}\) este independentă;
  \item \(M = N \oplus \bigl(\bigoplus_{j\in J} S_j\bigr)\).
\end{enumerate}
\end{proposition}

\begin{corollary}
Cu notaţiile de mai sus, pentru modulul semisimplu \(M\) există \(J \subseteq I\) astfel încât
familia \((S_j)_{j\in J}\) este independentă şi
\[
  M = \bigoplus_{i\in I} M_i.
\]
\end{corollary}

\begin{corollary}
Dacă \(M\) este un \(R\)-modul semisimplu şi \(N\) un submodul al său, atunci \(N\) şi
\(M/N\) sunt semisimple.
\end{corollary}

\begin{corollary}
O sumă directă de module semisimple este modul semisimplu.
\end{corollary}

\begin{theorem}
Fie \(M\) un \(R\)-modul. Următoarele afirmaţii sunt echivalente:
\begin{enumerate}
  \item \(M\) este semisimplu;
  \item \(M\) este izomorf cu o sumă directă de module simple;
  \item orice submodul al său este sumand direct în \(M\);
  \item orice şir exact
  \[
    0 \longrightarrow M' \xrightarrow{\,f\,} M \xrightarrow{\,g\,} M'' \longrightarrow 0
  \]
  este scindabil.
\end{enumerate}
\end{theorem}

\begin{definition}
Suma submodulelor simple ale lui \(M\) se numeşte \emph{soclul} lui \(M\) şi se notează
\(\soc(M)\). Dacă \(M\) nu conţine nici un submodul simplu atunci punem \(\soc(M) = 0\).
\end{definition}

\begin{proposition}
Fie \(M\) şi \(N\) două \(R\)-module şi \(f\colon M \to N\) un morfism.
Atunci \(f(\soc(M)) \subseteq \soc(N)\).
\end{proposition}

\begin{proposition}
Fie \(M\) un \(R\)-modul şi \(N\) un submodul al său. Atunci
\[
  \soc(N) = \soc(M) \cap N.
\]
\end{proposition}

\begin{proposition}
Dacă \(M = \bigoplus_{i\in I} M_i\), atunci
\[
  \soc(M) = \bigoplus_{i\in I} \soc(M_i).
\]
\end{proposition}

\begin{proposition}
Fie \(R\) un inel. Atunci \(\soc(R_R)\) este un ideal bilateral al lui \(R\).
\end{proposition}


% Capitolele următoare vor fi completate pe măsură ce avansăm.
% Exemplu:
% \chapter{Titlul capitolului 1}
% ...

\backmatter

% BIBLIOGRAFIE
% Când vom recrea lista de referinţe, putem folosi fie
% \begin{thebibliography}{99} ... \end{thebibliography}
% fie un fişier .bib separat.
% \bibliographystyle{plain}
% \bibliography{bibliography}

\end{document}
